\begin{table}[t]
\caption{Evaluation criteria and where they are measured.}
\label{tab:criteria}
\footnotesize
\begin{tabularx}{\columnwidth}{@{}>{\raggedright\arraybackslash}p{0.27\columnwidth}Lc@{}}
\toprule
Criterion & How it is measured & Section \\
\midrule
Functional correctness & Canonical scenarios, the incident sequence, a boundary matrix of malformed inputs, re-execution of every decision & 6.2, 6.3 \\
Tamper evidence & Attacks T1--T9 on the credit-scoring ledger, in each configuration & 6.4 \\
Cost & A four-arm ablation of the integrity layer; anchoring on a Fabric network & 6.5, 6.6 \\
Auditability & Fourteen audit questions (Q1--Q14) derived from the legal text, in four configurations & 6.7 \\
Controlled disclosure & Redaction tests; questions Q13 and Q14 on the redacted export & 5.5, 6.7 \\
Interoperability and conformance & Validation against the OSCAL schemas; parity between the Python issuer and the Go chaincode & 6.8 \\
\bottomrule
\end{tabularx}
\end{table}
